\renewcommand{\footerfilename}{ubuntu.tex}

\section{General Ubuntu Notes}


Canonical is company that creates, maintains, and offers support contracts for
the Ubuntu distrubution of Debian Linux. Called a 'distro', Canonical staff 
strive to make sure all packages are fully interoperable by matching revision
numbers for the many constitutient components. The images tend to lag the 
latest release levels. In addition to maintaining update integrity the distro
allows customers to better manage Current Vunerablities and Exposures (CVEs).

The steps with Ubuntu are:
\vspace{-.15cm}
\begin{enumerate}\addtolength{\itemsep}{-0.5\baselineskip}
%\setcounter{enumi}{N}
   \item   sudo apt update -- bring down the list of packages, but do not install anything.
     this helps to assure the package version numbers are all matched.
   \item   sudo atp upgrade -- bring (upgrade) the system's installed packages to the latest
     level consistent with the integration for the distro.
\end{enumerate}

The update step is fast, the upgrade step may take a while.

The changes are applied directly to the main distribution's file system. These are not
buffered. 

Make sure to backup the system to avoid any issues that may popup.

Reboot may be necessary.

\subsection{SNAP packages}

Snaps are relativly contained images that provide a degree of isolation between 
libraries that are not at the level of their host system. While not a virtual
box they do provide a degree of self-containment. 

\subsection{Linux Containeds (LXC}

Linux Containers come in two flavores: LXC and LXD. The LXD (Linux
Container Daemons) container is a 'lighter version' of the more
complete LXC environments.

\subsection{Docker}

A docker is a heaver container management system. They permit full install of
the OS together with their supported packages. They offer protection for
browser based applications -- an online hack destroys the container, but
a new image is easily (re)loaded and work may proceed apace. They also permit
using packages designed for other platforms. They are lighter than a full VirtualBox,
but have restrictions w.r.t. graphics and tedious access to their host systems' 
resources.

An example of using a Docker container is to download the Linux, Latex and
all supporting code to support Sphinx development without needing to have
that subsystem installed in a Python venv.





